\specialchapter{不变性表 --- I}

在本表及下表中，每一行对应于环可能具有的一个性质，每一列对应于由环 A 导出的一个环（p 表示 A 的一个素理想，S 表示 A 的一个不含 0 的乘性子集，而 A' 表示 A 在 A 的分式域 K 的有限代数扩张 L 中的整闭包）。假定环 A 具有该行所指的性质；本行与某列交叉处的 ``是''（相应地 ``否''、``？''）表示：由该列所指过程从 A 构造的每个环都具有该行所指的性质，这一事实为真（相应地为假、仍未知）。

参考文献指出本书或《代数》一书中证明所述结果的位置；对于下列两部著作中涉及正文或习题未提及的结果，情形同样如此：

\begin{enumerate}
\def\labelenumi{(\arabic{enumi})}
\item
  A. GROTHENDIECK, 代数几何原理，第 IV 章 (Publ. Inst. Htes Études Scient.，第 20 和 24 号，1964)。
\item
  M. NAGATA, 局部环，Interscience (New York)，1962。
\end{enumerate}

A/p

\(S^{-1}A\)

A{[}X{]}

A{[}{[}X{]}{]}

A'

A 为主理想整环

是

是

否 代数 VII, § 1, 习题 1

否

否 代数 VII, § 1, 习题 12

A 为 Dedekind 整环

是

是

否 代数 VII, § 1, 习题 1

否 VII, § 1, 习题 9

是 VII, § 2, 命题 9 的推论 2)

A 为因子分解整环

否 V, § 1, 习题 9

是 VII, § 3, 命题 3

是 VII, § 3, 定理 2

否 VII, § 3, 习题 9

否 代数 VII, § 1, 习题 12

A 为 Noether 整闭整环

否 V, § 1, 习题 9

是 V, § 1, 命题 16 的推论 1

是 V, § 1, 命题 13 的推论 1

是 V, § 1, 命题 14

？（若 L 可分则是；V, § 1, 命题 18 的推论 1）

A 为域或离散赋值环

是 VI, § 3, 第 6 节

是 VI, § 3, 第 6 节

否 代数 VII, § 1, 习题 1

否 VII, § 1, 习题 9

否 V, § 2, 习题 6

A 为域或高度 1 的赋值环

是 VI, § 4

是 VI, § 4, 命题 1

否 代数 VII, § 1, 习题 1

否 VII, § 1, 习题 9

否 V, § 2, 习题 6

A 为赋值环

是 VI, § 1, 定理 1

是 VI, § 1, 定理 1

否 代数 VII, § 1, 习题 1

否 VII, § 1, 习题 9

否 V, § 2, 习题 6

A 为完备赋值环

是 VI, § 5, 命题 1

是 VI, § 7, 命题 3

否 代数 VII, § 1, 习题 1

否 VII, § 1, 习题 9

否 V, § 2, 习题 6

A 为 Krull 整环

否 V, § 1, 习题 9

是 VII, § 1, 命题 6

是 VII, § 1, 命题 13

是 VII, § 1, 习题 9

是 VII, § 1, 命题 12

